% Proposed rewrite for §1.0 v2 L74 — "being-affected / different logos" sentence + footnote.
% Replaces from "The capacity for being-affected..." through the end of the footnote `}`.
% \hl{...} highlight removed (assumed WIP marker, no longer needed). Restore if intended for keeps.
% v2 (revised per feedback): heavy hyphenated-italicized labels in footnote rewritten as plain descriptive phrases.
% One LaTeX block on one line below (the actual sentence + footnote — paste as-is):

The capacity for being-affected is not a secondary feature added to cognition but a co-equal articulation of the very same actuality that cognition itself articulates, given under a different \textit{logos}: where cognition articulates the actuality as the soul's grasping of what is, being-affected articulates the same actuality as the soul's being-moved by what is. The two are not sequenced operations stacked on a neutral cognitive substrate but two \textit{logoi}---two structural accounts---of one indivisible moment in which the soul, as being-in-the-world, is at once cognitively grasping and affectively taken.\footnote{The claim that cognition and being-affected articulate the same underlying actuality under different \textit{logoi} draws together three Aristotelian moves that jointly secure the structural feature. (i) The self-applicative (or reflexive) character of soul-inquiry (\cite{aristotle2014soul} I.1, 402a4--10): the inquiry into the soul's essence is itself an exercise of the \textit{noetic} capacity it investigates, so that even the philosophical articulation of these structures is itself a specification under a \textit{logos} of the very actuality being investigated---a self-applicative point that prevents the analysis from treating cognition as if it were given from some external vantage. (ii) The co-given (or coinciding) actualization of perceiver and perceptible (\cite{aristotle2014soul} II.5, 425b26--426a26; developed at III.2): ``The actualization of the sensible object and that of the percipient sense is one and the same actualization, though their being is not the same'' (III.2, 425b26). The qualifier ``their being is not the same'' names precisely the structural possibility that one actuality admits of two distinct accounts: the same actualization, articulated as the perceptible's becoming actually perceived (the cognitive \textit{logos}), is at once articulated as the percipient's becoming actually affected (the affective \textit{logos}). (iii) The higher-order awareness of perception (\cite{aristotle2014soul} III.2, 425b12--25): ``since it is through sense that we are aware that we are seeing or hearing, it must be either by sight that one is aware of seeing, or by some sense other than sight.'' Awareness of perceiving is not a second operation appended to perception but a further articulation of the perceptual actuality itself---the soul's reflexive grasp of its own being-affected. Taken together, the three moves establish that cognition and being-affected are not sequenced operations stacked on a neutral cognitive substrate but two articulations (\textit{logoi}) of one structural moment in which the soul, as being-in-the-world, is at once cognitively grasping and affectively taken by what discloses itself to it---the precise point at issue in the anti-internalist thesis above.}
